\documentclass[12pt,a4paper]{article}
\usepackage[margin=1in]{geometry}
\usepackage{amsmath,amssymb,amsthm,booktabs,array}
\usepackage[hidelinks]{hyperref}
\hypersetup{pdftitle={Renormalization and Imprint Filtering: Source Identification, Memory, and Calibration},pdfauthor={Sidong Liu}}
\newtheorem{theorem}{Theorem}[section]
\newtheorem{proposition}[theorem]{Proposition}
\newtheorem{corollary}[theorem]{Corollary}
\newtheorem{definition}[theorem]{Definition}
\newcommand{\E}{\mathbb E}
\newcommand{\Var}{\operatorname{Var}}
\title{Renormalization and Imprint Filtering:\\
Source Identification, Memory, and Calibration}
\author{Sidong Liu, PhD\\iBioStratix Ltd\\
\texttt{sidongliu@hotmail.com}}
\date{October 2026}
\begin{document}
\emergencystretch=2em
\raggedbottom
\maketitle

\begin{abstract}
A physical trace is readable relative to a source question and a
permitted measurement. We make that relation operational on a
continuing lattice field: bounded randomized forces at two durations
separate an unknown source relaxation rate from an unknown positive
reader gain, and the acquired response history selects the next read
duration and physical force. The changed carrier continues through
calibration and action. A finite martingale certificate retains its
nonstationary contribution and tolerates unknown centered independent
read noise with a known variance cap.
Exact Gaussian block spectra establish the scale-dependent distinctions
on which this inquiry acts. Coarse readings can retain mass information
while suppressing higher-gradient contrasts, and mixing relaxation rates
makes retained history useful for prediction. Independent preparations
give separate gain-free decision laws and explicit low-gain budget
obstructions. History also changes which reading is acquired: five
retained bits beat every fixed allocation in a slowly changing source
family. On one field trajectory, a paid block reading or no new reading
is followed by a bounded force. Gaussian conditioning gives the optimal
one-round rule, and sixteen-bit records preserve its strict advantage
over every history-independent reading choice in the stated case.
The noise extension concerns response calibration; the Gaussian
likelihood and conditional-state results keep their original noise
assumptions. Reading prices, mechanical switching work and the supplied
apparatus remain separate parts of the resource account.
\end{abstract}

\section{Which traces survive a change of scale?}

The world can carry a history without presenting that history in a form
an observer can use. Large-scale fluctuations may identify a material's
phase while revealing little about its microscopic couplings. A present
measurement may be insufficient to predict the next one even when the
physical dynamics is Markovian. A changing signal may reflect either
the source or the instrument. These are different failures of
readability.

The imprint proposal asks how generative conditions remain legible in
the structures they produce \cite{WorldImprint}. Renormalization makes
scale central to this question: an effective description preserves
selected distinctions and suppresses others. Information geometry
already relates coarse-graining to parameter distinguishability
\cite{Machta2013,Raju2018}. Projection methods likewise establish that
eliminating degrees of freedom can introduce memory \cite{Mori1965}.
We use these ideas together, with the reading experiment and its
resources specified.

Our central claim is that an imprint is a relation among a source
contrast, its surviving carrier, and a reading interface. It is not an
additional substance carried by every organized object. A single
lattice model makes the relation calculable. Its static distribution
gives source-identification tests; its dynamics gives a value for
retained history; an uncertain gain gives a calibration problem.
Information can survive physically without being readable through a
given interface, and successful world-reading can require a model of
the reader's own coupling.

Block averaging is an exact transformation of the Gaussian probability
law, not a claim that the effective action stays within a short list of
local couplings. The mathematical tools are standard. Their common
implementation establishes what scale filtering, memory, and
calibration jointly imply for the imprint account. In particular,
records can serve as references for one another through either spatial
scale or controlled readout settings. Whether this removes an unknown
gain depends on the actual record distribution, not just on a cost
assigned to each reading.

The initial calculations use three experiments in the same model family:
independent static preparations, the subsequent relaxation of a
preparation in a common bath, and prediction along a stationary
trajectory. Their inputs are stated separately. They are not claimed
to be simultaneous observations of one run.
Section~\ref{sec:history-scale} adds a fourth experiment: a persistent
regime changes which of two scales carries the current source contrast.
Its retained history changes the next allocation of real readings.
This is not information recovered by postprocessing a discarded field.
Section~\ref{sec:field-control} instead follows one field through past
observation, a selected current reading and a subsequent force. Its
control changes the carrier that will be read again.
Section~\ref{sec:active-reference} uses that actual change as a
calibration reference. It separates source dynamics from reader gain
on one spectral mode, then uses the response history to choose a
future temporal scale and control. These are declared protocols on
one lattice family, not simultaneous free access to every observable.

\section{A source family and its block observations}

On $\Lambda=(\mathbb Z/L\mathbb Z)^2$, take a real centered Gaussian
field with unitary Fourier convention
$\widehat X(k)=L^{-1}\sum_x e^{-ik\cdot x}X_x$ and precision
\begin{align}
 K_\theta(k)&=\mu+s(k)+g\,s(k)^2+\gamma\,a(k),\label{eq:K}\\
 s(k)&=4\sin^2(k_1/2)+4\sin^2(k_2/2),\nonumber\\
 a(k)&=16\sin^4(k_1/2)+16\sin^4(k_2/2).\nonumber
\end{align}
Here $\mu>0$, $g,\gamma\geq0$ suffice for positivity. Derivatives at
$g=\gamma=0$ use the open finite-lattice extension where $K_\theta>0$.
No claim about efficient estimation at a parameter-space boundary is
needed. The parameters describe preparation conditions, not every
detail of a sample's formation history.

The coefficients $g$ and $\gamma$ change higher-gradient terms.
For example, $\gamma=1/12$ cancels the fourth-order anisotropic
small-momentum error of the nearest-neighbor Laplacian.

For $\ell\mid L$, observe nonoverlapping means
$Y_j=\ell^{-2}\sum_{a_1,a_2=0}^{\ell-1}X_{\ell j+a}$.
There are $N_\ell=(L/\ell)^2$ coarse coordinates. For coarse momentum
$Q$, its aliases are $k_n=(Q+2\pi n)/\ell$, with
$n\in\{0,\ldots,\ell-1\}^2$. Set
$W_\ell(k)=\sum_{a=0}^{\ell-1}e^{ika}$.

\begin{proposition}[Exact block spectrum]\label{prop:block}
The unitary Fourier covariance of the coarse field has eigenvalues
\begin{equation}
 v_{\theta,\ell}(Q)=\ell^{-6}\sum_n
 \frac{|W_\ell(k_{n,1})W_\ell(k_{n,2})|^2}{K_\theta(k_n)}.
 \label{eq:block}
\end{equation}
\end{proposition}
\begin{proof}
Insert the inverse fine Fourier transform in the block mean and take
the coarse unitary transform. The coarse-site sum selects the alias
set. Each fine coefficient is multiplied by
$\ell^{-3}W_\ell(k_{n,1})W_\ell(k_{n,2})$. Squaring and summing its
variance gives \eqref{eq:block}.
\end{proof}

For white noise $K=1$, this yields $v_{\theta,\ell}=\ell^{-2}$.
The zero coarse Fourier mode has variance $1/(\ell^2\mu)$; the literal
global spatial mean has variance $1/(L^2\mu)$. Both are blind to $g$
and $\gamma$. Nontrivial fluctuations can therefore be completely
uninformative about a specified source contrast.

Fourier coefficients of a real field occur in conjugate pairs.
Information sums below count the corresponding independent real sine
and cosine coordinates, not independent complex observations twice.
The effective precision is the inverse coarse covariance. It generally
contains terms beyond \eqref{eq:K}; we do not assume a closed
three-parameter RG flow.

\section{Information retained by a coarse description}

\begin{definition}[Readable source contrast]
A source contrast is readable through a fixed protocol if its two
preparations induce different record distributions. Its strength can be
expressed by decision risk or divergence. A resource claim also states
how the records are acquired, retained, and processed.
\end{definition}

\begin{proposition}[Gaussian information]\label{prop:information}
One complete coarse-field snapshot has Fisher information
\begin{equation}
 F_{ij}^{(\ell)}=\frac12\sum_Q
 \partial_i\log v_{\theta,\ell}(Q)\,
 \partial_j\log v_{\theta,\ell}(Q).
 \label{eq:F}
\end{equation}
For two preparations with spectra $v_0,v_1$, its Chernoff information is
\begin{equation}
 C_\ell=\max_{0\leq u\leq1}\frac12\sum_Q
 \left[u\log v_0+(1-u)\log v_1+
 \log\left(\frac{u}{v_0}+\frac{1-u}{v_1}\right)\right].
 \label{eq:C}
\end{equation}
It is zero exactly when the two coarse covariances coincide.
\end{proposition}
\begin{proof}
In real Gaussian coordinates the variance score is
$\frac12\partial_i\log v\,(Z^2/v-1)$; its second moment gives
\eqref{eq:F}. Integrating $p_0^u p_1^{1-u}$ gives \eqref{eq:C}.
Equality at zero requires equal Gaussian laws.
\end{proof}

A zero Fisher direction at one parameter value means only first-order
local indistinguishability. It need not describe a globally unobservable
source family. The zero mode above, by contrast, is exactly blind to all
$g,\gamma$ contrasts at fixed $\mu$.

Nested block averages of the same sample satisfy
\[
 F^{(\ell_2)}\preceq F^{(\ell_1)},\qquad
 C_{\ell_2}\leq C_{\ell_1}\quad(\ell_1\mid\ell_2).
\]
The coarse score is the conditional expectation of the finer score;
conditional variance proves the first inequality. Conditional
H\"older inequality proves the second. Known parameter-independent
coordinate or amplitude rescaling changes neither information.
Acquiring a different observable or preparing more fields changes the
experiment and is not postprocessing of the old records.

\section{Reading cost and spectral access}

Spatial block values are correlated. Thus $F^{(\ell)}/N_\ell$ and
$C_\ell/N_\ell$ are densities for a complete snapshot, not information
in arbitrary independent block readings. For $n$ independently
prepared snapshots, the equal-prior optimal error obeys
\begin{equation}
 P_e\leq\tfrac12 e^{-nC_\ell}.
 \label{eq:Pe}
\end{equation}
A $5\%$ sufficient certificate needs
$\lceil\log(10)/C_\ell\rceil$ snapshots and $N_\ell$ times as many scalar
records. Correlated snapshots from one trajectory cannot replace
independent preparations.

Table~\ref{tab:scale} uses $L=256$, $\theta_0=(.01,0,0)$.
The mass contrast changes $\mu$ to $.0125$; the microscopic contrast
changes $\gamma$ to $1/12$. This table idealizes continuous precision.
\begin{table}[ht]
\centering\small
\begin{tabular}{rrrrrr}
\toprule
$\ell$&$N_\ell$&$F_{\mu\mu}$&$F_{\gamma\gamma}$&
$C_{\rm mass}$&$C_{\rm micro}$\\
\midrule
1&65536&262732.0&248831.1&.183030&170.010\\
4&4096&256284.8&433.963&.178014&.333084\\
8&1024&242440.5&8.50597&.167361&.00671261\\
16&256&202231.1&.133228&.137232&.000107754\\
32&64&126311.3&.00206272&.0831354&.00000168534\\
\bottomrule
\end{tabular}
\caption{Information in one complete block snapshot.}\label{tab:scale}
\end{table}

For the mass contrast, changing $\ell=1$ to $32$ reduces the sufficient
record count from $851968$ to $1792$, while increasing preparations
from $13$ to $28$. For the microscopic contrast, preparations grow from
$1$ to $1366242$. The former is a trade between preparation and readout,
not information gained by discarding data. Preparation count, physical
support, reads, precision, retention, and computation belong in
separate resource entries.

Changing spectral access gives another comparison. For the microscopic
contrast, a real coordinate at $k=(\pi,0)$ has variances $.24937656$
and $.18714910$. Use $255$ fixed thresholds equally spaced in
$[-8\sqrt{v_{\max}},8\sqrt{v_{\max}}]$, giving $256$ outcomes including
the two unbounded tails. Integrating the Gaussian laws over these bins
gives Chernoff information $.00514034$ per independent reading.
Then $448$ independent preparations suffice for the $5\%$ certificate,
with $3584$ bits of recorded outcomes. Counters, decoding, preparation,
and the full spatial support of the spectral observable are additional
resources. One recorded number is not a single-site interaction.

The low coordinate $k=(2\pi/256,0)$ instead has continuous Chernoff
information about $5.08\times10^{-13}$; the zero mode has none.
A new interface reveals a distinction by accessing the carrier again,
not by reconstructing discarded data from the old coarse records.

\section{Reading a past preparation after relaxation}

Prepare $X_0\sim N(0,K_\theta^{-1})$ and subsequently use a common bath
with fixed $K_0=K_{\theta_0}$, independent of the preparation parameter
$\theta$:
\begin{equation}
 dX_t=-K_0X_t\,dt+\sqrt2\,dW_t.
 \label{eq:OU}
\end{equation}
Physical loss of preparation history and spatial coarse-graining are
now distinct operations in the same process.

\begin{proposition}[Remaining preparation information]
The fine variance at time $t$ is
\begin{equation}
 V_\theta(k,t)=K_0(k)^{-1}+e^{-2K_0(k)t}
 [K_\theta(k)^{-1}-K_0(k)^{-1}].
 \label{eq:history}
\end{equation}
Replace $K_\theta^{-1}$ by $V_\theta$ in \eqref{eq:block} for its block
spectrum. At $\theta=\theta_0$, its log-variance derivative is
\begin{equation}
 r_i(Q,t)=-
 \frac{\sum_n |W_1W_2|^2e^{-2K_0t}(\partial_iK_\theta)/K_0^2}
 {\sum_n |W_1W_2|^2/K_0}.
 \label{eq:score}
\end{equation}
\end{proposition}
\begin{proof}
The solution is
$e^{-K_0t}X_0+\sqrt2\int_0^t e^{-K_0(t-s)}dW_s$.
Its independent covariance contributions give \eqref{eq:history}.
Differentiate and project to obtain \eqref{eq:score}.
\end{proof}

The source parameter is no longer being applied after preparation.
For $K_0(k)\geq\mu>0$, its full-field Fisher matrix at the reference
satisfies
\[
 F_{\rm full}(t)\preceq e^{-4\mu t}F_{\rm full}(0).
\]
Each mode contributes a positive rank-one matrix multiplied by
$e^{-4K_0(k)t}$. A fixed reader adds further loss; this formula does
not compare every coarse reader against its own initial information.

In the reference field of Table~\ref{tab:scale}, full-field $\gamma$
information falls from $248831.1$ at $t=0$ to $49.9307$ at $t=1$
and $.0237407$ at $t=10$. Mass information at $t=10$ remains
$101633.8$. A readable history has a depth relative to its source
question, not a single universal lifetime. The bath and elapsed time
are known inputs.

\section{The prediction value of a retained record}

Write $K=K_0$ and start \eqref{eq:OU} in its stationary law $N(0,K^{-1})$.
For a real coarse Fourier coordinate $Y(t)$ of variance $v>0$, the
normalized correlation is
\begin{equation}
 r(t)=\sum_n p_n e^{-K(k_n)|t|},\qquad
 p_n=\frac{|W_1W_2|^2/K(k_n)}
 {\sum_m|W_1W_2|^2/K(k_m)}.
 \label{eq:r}
\end{equation}

\begin{theorem}[Memory relative to the coarse interface]\label{thm:memory}
Fix $\Delta>0$, write $Y_j=Y(j\Delta)$, $r_1=r(\Delta)$ and
$r_2=r(2\Delta)$. Then
\begin{equation}
 d=r_2-r_1^2=\Var_p(e^{-K\Delta})\geq0.
 \label{eq:d}
\end{equation}
The minimum square risks for predicting $Y_1$ from $Y_0$ and from
$(Y_0,Y_{-1})$, over all measurable predictors, are
\begin{equation}
 R_1=v(1-r_1^2),\qquad
 R_2=R_1-\frac{v d^2}{1-r_1^2}.
 \label{eq:risks}
\end{equation}
The improvement is strict exactly when the positive-weight alias set
has more than one relaxation rate.
\end{theorem}
\begin{proof}
Equation~\eqref{eq:d} follows from the positive mixture \eqref{eq:r}.
Positivity of $K$ and $\Delta$ gives $0<r_1<1$.
The three records are jointly Gaussian with correlations $r_1,r_1,r_2$.
Their conditional expectation minimizes square loss among all measurable
functions. Its Schur-complement variance gives \eqref{eq:risks};
the minimizing two-record predictor is
\begin{equation}
 \widehat Y_1=a_1Y_0+a_2Y_{-1},\qquad
 a_1=\frac{r_1(1-r_2)}{1-r_1^2},\quad
 a_2=\frac{d}{1-r_1^2}.
 \label{eq:pred}
\end{equation}
A mixture variance vanishes exactly when the rates on its support agree.
\end{proof}

This also characterizes Markov closure. Distinct coarse real Fourier
coordinates are independent processes. If each alias group has one
rate, every coordinate is an Ornstein--Uhlenbeck process. Otherwise
\eqref{eq:d} gives conditional dependence of past and future after
conditioning on the entire current coarse field, so those variables
alone are not Markov. The complete microscopic state remains Markov.

For $L=64$, $\ell=4$, $\mu=.1$, $g=\gamma=0$, coarse Fourier index
$(3,5)$ and $\Delta=.2$, the results are
\[
 v=.10868306,\quad d=.005993705,\quad
 R_1=.0214457374,\quad R_2=.0214259507.
\]
The improvement is $.0923\%$.

\subsection{Finite storage and arithmetic}

\begin{corollary}[Finite record precision]\label{cor:finite}
Quantize each record to $2^B$ uniform levels in $[-A,A]$, clipping
outside. Put
\[
 \delta=\frac{2A}{2^B-1},\quad z=A/\sqrt v,\quad
 \eta=\delta/2+\sqrt{2v[z\phi(z)+\overline\Phi(z)]},
\]
where $\phi,\overline\Phi$ are the standard normal density and upper
tail. Applying \eqref{eq:pred} to these records has risk
\begin{equation}
 \widetilde R_2\leq R_2+(|a_1|+|a_2|)^2\eta^2.
 \label{eq:finite}
\end{equation}
If the right side is below $R_1$, this finite-record predictor beats
every predictor using the exact current record alone.
\end{corollary}
\begin{proof}
The quantization error has $L^2$ norm at most $\eta$, by the Gaussian
tail identity
$\E[Y^2\mathbf1_{\{|Y|>A\}}]=2v[z\phi(z)+\overline\Phi(z)]$.
The conditional-prediction residual is independent of both records and
hence orthogonal to their quantization errors. Its risk adds to the
predictor perturbation's second moment. The $L^2$ triangle inequality
then gives \eqref{eq:finite}.
\end{proof}

In this example, $B=10$ and $A=8\sqrt v$ give
$\eta\leq .00257817$ and $\widetilde R_2\leq .021431324<R_1$.
Only the retained data occupy $20$ bits. The full-support reading
operation, model coefficients, and computation still have costs.

The rational-arithmetic construction in Appendix~\ref{app:force-precision}
retains this strict prediction advantage with finite coefficient and
output registers.

The mechanism is different from asserting that every viable system
needs memory. Coarse observation removes the sufficiency of the
present. A past record then supplies information missing from that
present, and its predictive value can be computed.

\section{Source identification and instrument calibration}

The previous calculations assumed a known reading interface. Now take
$K(k)=\mu+s(k)$ and an unknown fixed positive instrument gain $b$.
A real spectral reading has law
\[
 Z_i\sim N\!\left(0,\frac{b^2}{\mu+\lambda_i}\right),
 \qquad \lambda_i=s(k_i).
\]
The same gain acts at both known frequencies. A positive sign is
stipulated because centered variance data cannot determine it.

\begin{proposition}[Source--gain separation]\label{prop:gain}
One spectral band cannot identify $\mu,b$ separately.
Two bands with $\lambda_1\ne\lambda_2$ identify both globally.
For $n_1,n_2>0$ independent preparations, let $\eta_b=\log b$ and
$q_i=-1/(\mu+\lambda_i)$. The Fisher matrix and its Schur complement are
\begin{align}
 F&=\frac12\sum_i n_i
 \begin{pmatrix}q_i\\2\end{pmatrix}\begin{pmatrix}q_i&2\end{pmatrix},\\
 F_{\mu\mid\eta_b}
 &=\frac{n_1n_2}{2(n_1+n_2)}(q_1-q_2)^2>0.
 \label{eq:gain}
\end{align}
At a fixed total reading count the latter is maximized by equal
allocation.
\end{proposition}
\begin{proof}
One variance determines only $b^2/(\mu+\lambda_i)$.
For two population variances set $\rho=v_1/v_2$. Then
\[
 \rho=\frac{\mu+\lambda_2}{\mu+\lambda_1},\qquad
 \mu=\frac{\lambda_2-\rho\lambda_1}{\rho-1},\qquad
 b^2=v_1(\mu+\lambda_1).
\]
This determines the positive gain uniquely. Apply \eqref{eq:F} to
the log variances and take a Schur complement to obtain \eqref{eq:gain}.
For fixed $n_1+n_2$, their product is maximal at equal allocation.
\end{proof}

For $L=64$, $\mu=.1$, modes $(10\pi/64,0)$ and $(\pi/2,0)$,
and $64$ preparations per band, $F_{\mu\mid\eta_b}=99.88849656$.
Population identifiability and finite-sample information do not assert
that a sample-variance ratio is unbiased or attains an information bound.

The spectral family, band locations, gain stability, and common
frequency response are inputs. Their importance is physical:
a system can distinguish source change from gain change because these
alternatives leave different joint records. To exercise that capacity,
it must estimate, retain the result, and revise subsequent reading or
action. T-DOME~III and Foundation~III provide separate drift and
feedback constructions for those operations \cite{TDIII,FoundationIII}.
The existence of an identifiable model is not itself an agent that
uses it.

The continuing experiment in Section~\ref{sec:active-reference} turns
that distinction into the next operation. A fresh force coin supplies
an intervention reference on a mode whose state is retained. Responses
at two durations identify its source rate and reading gain; the
inferred values choose a new time scale and force on the field already
changed by calibration. The same finite certificate allows unknown
centered read noise with a known variance cap. This is an additional
response-based route, rather than a use of the independent-preparation
variance laws on correlated readings.

\subsection{Block scales as references for one another}
\label{subsec:two-scales}

Return to the three preparations of Table~\ref{tab:scale}, denoted
$H_0$, $H_\mu$, and $H_\gamma$. One reading now means one scalar
block mean from a freshly prepared field, rather than a complete
coarse-field snapshot. Translation invariance gives its variance
\[
 \sigma_\ell^2(H)=\frac1{N_\ell}\sum_Q v_{H,\ell}(Q).
\]
With a common unknown gain $b>0$, the record has law
$Y\sim N(0,b^2\sigma_\ell^2(H))$. All preparations are independent;
the mean is known to be zero. Fix the reading counts
$n_\ell,n_m>0$ in advance and let $S_\ell,S_m$ be the sums of squares.

\begin{proposition}[Gain-invariant source decisions]\label{prop:ratio}
For every $b>0$,
\begin{equation}
 R=\frac{S_\ell/n_\ell}{S_m/n_m}
 \ \overset d=\ \rho_H F_{n_\ell,n_m},\qquad
 \rho_H=\frac{\sigma_\ell^2(H)}{\sigma_m^2(H)}.
 \label{eq:ratio}
\end{equation}
Every function of the two positive sufficient statistics that is
invariant under common positive rescaling depends only on $R$.
With one such statistic, every invariant function is constant.
Distinct $\rho_H$ identify finitely many sources consistently as
both reading counts tend to infinity, without a prior on $b$.
\end{proposition}
\begin{proof}
The variables $S_\ell/[b^2\sigma_\ell^2(H)]$ and
$S_m/[b^2\sigma_m^2(H)]$ are independent chi-squares. Dividing each
by its degrees of freedom gives the standard $F$ law
\cite{NISTF}. Positive rescaling is transitive on the positive reals;
on positive pairs its orbits are labeled by the ratio. Finally,
the two sample second moments converge to their population values,
so $R\to\rho_H$ almost surely.
\end{proof}

For equal source weights, choose the largest scaled $F$ density.
This rule minimizes the average error among gain-invariant rules
based on these sufficient statistics. When $n_\ell=n_m=n$, adjacent
ordered ratios $\rho_i<\rho_{i+1}$ have decision boundary
$\sqrt{\rho_i\rho_{i+1}}$: the scaled density is proportional to
$\rho^{n/2}R^{n/2-1}/(\rho+R)^n$.
Thus ordinary $F$ cumulative probabilities give each source's exact
finite-sample success probability.

At scales $(1,32)$ the ratios for $(H_0,H_\mu,H_\gamma)$ are
$12.85474928$, $14.55261454$, and $11.92780351$.
With $20{,}000$ readings split equally, the success probabilities are
$.968359$, $.999036$, and $.969323$. Table~\ref{tab:ratio-counts}
gives the smallest even total for which all three reach $.95$ under
the equal-split rule. These are numerical evaluations of the exact
decision law, not Monte Carlo estimates or minima over all allocations.

\begin{table}[ht]
\centering\small
\begin{tabular}{rr}
\toprule
Scale pair & Total independent readings\\
\midrule
$(1,32)$ & $15{,}984$\\
$(1,16)$ & $20{,}082$\\
$(1,8)$ & $36{,}732$\\
$(1,4)$ & $107{,}746$\\
$(2,32)$ & $96{,}336$\\
$(4,32)$ & $758{,}974$\\
\bottomrule
\end{tabular}
\caption{Equal-split reading counts for at least $.95$ success under
each source. Each scalar record requires a new field preparation;
gain is fixed but unknown, and readout is noiseless.}
\label{tab:ratio-counts}
\end{table}

The ratio removes gain from the decision; it does not make gain
unidentifiable from the full data. Once the source is identified,
its absolute variance determines $b$. Nor is $R$ an unbiased estimate
of $\rho_H$: its mean is
$\rho_H n_m/(n_m-2)$ when $n_m>2$.
The decision uses the distribution, not that mean.
Preparation, block support, record precision, and subsequent storage
remain separate resources. The fixed-count law cannot be reused
unchanged after data-dependent stopping or allocation.

\subsection{Readout noise and the price of a reference}
\label{subsec:noise}

A noisy probe changes both the cost and the distribution of a record.
Consider independent Gaussian readout noise $\xi\sim N(0,r)$,
with known variance. Two possible placements give
\begin{align}
 Y=b(Z+\xi):&\quad V=b^2(\sigma_\ell^2(H)+r),\label{eq:noise-in}\\
 Y=bZ+\xi:&\quad V=b^2\sigma_\ell^2(H)+r.\label{eq:noise-out}
\end{align}
These are alternative reading models. Take $r(E)=\nu/E$ with
$\nu>0$, and charge $E>0$ per reading. This energy--noise relation
is a specified probe model; preparation, memory, switching and the
rest of the apparatus are not included in this charge.

\begin{proposition}[Two kinds of calibration reference]\label{prop:noise}
For fixed settings in \eqref{eq:noise-in}, the gain-invariant ratio
law holds with
$\rho_H=[\sigma_\ell^2(H)+r_\ell]/[\sigma_m^2(H)+r_m]$.
These ratios must be distinct to identify the sources by this rule.
Moreover, at a single spatial scale, two known noise variances
$r_1\ne r_2$ identify both the gain and the source variance:
\begin{equation}
 b^2=\frac{V_1-V_2}{r_1-r_2},\qquad
 \sigma_\ell^2(H)=V_1/b^2-r_1.
 \label{eq:noise-reference}
\end{equation}
This identifies the source whenever its variances at that scale differ.
The ratio law also applies to the two settings at one scale: with
$\ell=m$ and $r_1\ne r_2$,
$\rho_H=[\sigma_\ell^2(H)+r_1]/[\sigma_\ell^2(H)+r_2]$ is strictly
monotone in $\sigma_\ell^2(H)$, so this reference also gives an exact,
gain-free finite-sample decision.

For fixed settings in \eqref{eq:noise-out}, the raw sample ratio has
law $(V_\ell/V_m)F_{n_\ell,n_m}$, generally dependent on $b$.
Nevertheless, for every fixed $b>0$,
\begin{equation}
 \frac{V_\ell-r_\ell}{V_m-r_m}
 =\frac{\sigma_\ell^2(H)}{\sigma_m^2(H)}.
 \label{eq:noise-subtracted}
\end{equation}
Distinct source ratios therefore still identify the source and then
the gain. Changing only the noise variance at a single scale in this
second model identifies $b^2\sigma_\ell^2(H)$, not its two factors.
\end{proposition}
\begin{proof}
The chi-square argument applies to either pair of variances, including
two settings at the same scale; there the derivative of
$(\sigma^2+r_1)/(\sigma^2+r_2)$ in $\sigma^2$ is
$(r_2-r_1)/(\sigma^2+r_2)^2\ne0$. In
\eqref{eq:noise-in}, subtracting the population variances at the two
settings gives \eqref{eq:noise-reference}; in
\eqref{eq:noise-out}, subtracting the known noise gives
\eqref{eq:noise-subtracted}. At a fixed scale in the second model,
every setting has variance $c+r$ with the same unknown product $c$.
All sources of positive variance can realize that product by changing
$b$.
\end{proof}

The subtractions in \eqref{eq:noise-subtracted} are statements about
population variances. Finite sample estimates can fall below the
noise floor and their subtracted ratio is not $F$ distributed.
Consistency at any fixed $b>0$ follows from convergence of sample
second moments; uniform accuracy down to arbitrarily small gains
requires a separate resource analysis. A proper prior on gain can
also give a single setting source information through absolute
amplitude, despite the unknown-fixed-gain ambiguity.

\begin{proposition}[Low-gain limit at fixed probe budget]\label{prop:budget}
Use \eqref{eq:noise-out} with $r(E)=\nu/E$ and two equally weighted
sources $H,H'$. Fix the same $b>0$ under both sources and a finite
set of allowed scales. At each of at most $T<\infty$ readings, a
policy $\pi$ may choose the scale and energy using previous records
and source-independent randomization, with
$\sum_i E_i\leq B$ on every history. Each reading uses an independent
new preparation, and no other source-dependent data are supplied.
Put $\delta=\max_\ell|\sigma_\ell^2(H)-\sigma_\ell^2(H')|$.
For the resulting transcript laws,
\begin{equation}
 D_{\rm KL}(P_H^\pi\Vert P_{H'}^\pi)
 \leq\frac{b^4\delta^2B^2}{4\nu^2},\qquad
 P_{\rm err}^{\pi}\geq
 \max\!\left\{0,\frac12-\frac{b^2\delta B}{4\sqrt2\,\nu}\right\}.
 \label{eq:budget}
\end{equation}
The error bound applies to every terminal decision, including
decisions based on adaptively chosen reads.
\end{proposition}
\begin{proof}
For centered Gaussian variances $v,w$,
\[
 D_{\rm KL}(N(0,v)\Vert N(0,w))
 =\tfrac12[v/w-1-\log(v/w)]
 \leq\frac{(v-w)^2}{4\min(v,w)^2}.
\]
This follows by Taylor's theorem for $-\log x$ between $1$ and $v/w$.
Conditional on a common transcript history, the policy uses the same
distribution of settings under either source. Each conditional
divergence is at most $b^4\delta^2E_i^2/(4\nu^2)$, because
$\min(v,w)\geq\nu/E_i$. The divergence chain rule and
$\sum_i E_i^2\leq B^2$ give the first bound. Stopping can be padded
with source-independent null records up to $T$. The equal-prior
testing error is at least $(1-\|P_H^\pi-P_{H'}^\pi\|_{\rm TV})/2$;
Pinsker's inequality gives the second bound.
\end{proof}

The bound uses only $\sum_iE_i^2\leq B^2$, which is tight when the
budget is spent on very few readings. If every reading also obeys
$E_i\leq E_{\max}$, then $\sum_iE_i^2\leq BE_{\max}$ and $B^2$ in
\eqref{eq:budget} may be replaced by $BE_{\max}$; keeping the pathwise
sum $\sum_iE_i^2$ gives a sharper statement still.

No allocation of this fixed probe budget, adaptive or otherwise,
maintains reliable discrimination uniformly as $b\to0$.
The obstruction belongs to the placement \eqref{eq:noise-out}. With
the noise before the gain, \eqref{eq:noise-in}, the same argument gives
$|v-w|\leq b^2\delta$ and $\min(v,w)\geq b^2\nu/E_i$, so the factor
$b^4$ cancels. Fixed designs with gain-invariant decisions, and
policies whose choices and decisions use only scale-invariant
statistics, then perform identically at every gain; a policy that
reacts to absolute amplitudes can still behave differently at
different gains.
This finite-resource obstruction coexists with identifiability at
each fixed positive gain. Conversely, \eqref{eq:noise-reference}
shows that a controlled change in the probe can itself provide a
reference. Cross-scale access is one way to resolve source--reader
ambiguity; a known change in the reading interaction is another.

\section{Self-reading at a chosen scale}\label{sec:self-scale}

Prediction and passive self-identification have the same statistical
quantity when the competing interface produces independent records
with the same one-record law. For a stationary record law $P$ and
that independent law $Q^{\otimes n}$, the divergence chain rule gives
\begin{equation}\label{eq:predictive-rate}
 D(P^n\Vert Q^{\otimes n})=
 nD(P_1\Vert Q)+\sum_{j=2}^n I_P(Y_j;Y_1,\ldots,Y_{j-1}).
\end{equation}
When $P_1=Q$, the divergence per record tends to
$I_P(Y_0;Y_{-\infty}^{-1})$, possibly infinite. This is the passive
test in \emph{The World as Imprint}~\cite{WorldImprint}, not an upper
bound on all active identification protocols. It also gives
Theorem~\ref{thm:memory} a self-reading interpretation: the increase
from one past Gaussian record to two is
$\frac12\log(R_1/R_2)>0$ exactly when the alias rates differ.

Now sample one physical block mean along a stationary trajectory of
\eqref{eq:OU}, at interval $\Delta$, and add independent Gaussian
readout noise of variance $r\geq0$. Write
\[
 w_\ell(k)=\frac{|W_\ell(k_1)W_\ell(k_2)|^2}
 {L^2\ell^4K(k)},\qquad a_k=e^{-K(k)\Delta},\qquad
 v_\ell=\sum_k w_\ell(k).
\]
The record spectral density is
\begin{equation}\label{eq:reading-density}
 f_{\ell,r}(\omega)=r+\sum_k w_\ell(k)
 \frac{1-a_k^2}{1-2a_k\cos\omega+a_k^2}.
\end{equation}
The positive finite mixture has absolutely summable correlations and
a strictly positive smooth density. The Gaussian entropy-rate formula
therefore applies~\cite{Gray2006}.

\begin{proposition}[The reader's noise changes its informative scale]
\label{prop:self-scale}
Against independent Gaussian records of variance $v_\ell+r$, the
passive divergence rate, in nats per continuous record, is
\begin{equation}\label{eq:scale-information}
 \mathcal I_\ell(r)=\frac12\log(v_\ell+r)
 -\frac1{4\pi}\int_{-\pi}^{\pi}\log f_{\ell,r}(\omega)\,d\omega.
\end{equation}
It is nonincreasing in $r$ at each scale. The scale maximizing it need
not remain fixed when the reader's noise changes.
\end{proposition}
\begin{proof}
The independent law has entropy rate
$\frac12\log[2\pi e(v_\ell+r)]$, and its cross-entropy under the
matched law is the same because their one-record marginals agree.
Subtract the matched Gaussian entropy rate to get
\eqref{eq:scale-information}. Alternatively this is the limiting
predictive information in \eqref{eq:predictive-rate}.
Adding independent noise of variance $r_2-r_1$ to every record maps
both candidate laws at $r_1$ to those at $r_2$. Data processing for
each finite transcript and passage to the rate prove monotonicity.
The table below supplies the change of maximizing scale.
\end{proof}

\begin{table}[ht]
\centering\small
\begin{tabular}{rrrrrrrr}
\toprule
$r$ & $\ell=1$ & $2$ & $4$ & $8$ & $16$ & $32$ & $64$\\
\midrule
$0$ & .3405 & .6219 & .9486 & 1.2351 & 1.4254 & 1.5250 & 1.6194\\
$10^{-4}$ & .3404 & .6210 & .9449 & 1.2200 & 1.3684 & 1.3390 & 1.1516\\
$10^{-2}$ & .3229 & .5479 & .7092 & .6769 & .4599 & .2076 & .0688\\
\bottomrule
\end{tabular}
\caption{Passive self-reading rate for $L=64$, $\mu=.1$, $g=\gamma=0$,
$\Delta=.2$. Each row compares fixed ideal continuous-record
interfaces, not a finite-memory controller or a total-cost optimum.}
\label{tab:self-scale}
\end{table}

The best listed scale moves from $64$ to $16$ to $4$. Coarse means
have longer temporal correlations, but smaller variance; fixed
readout noise masks them first. After noise rises from $10^{-4}$ to
$10^{-2}$, staying at scale $16$ gives $.4599$ nats per record,
whereas scale $4$ gives $.7092$. These numbers assume known noise.
They do not by themselves identify the damage, choose a scale online,
or give the support and apparatus bill for the new interface.

\section{A retained history chooses the next reading scale}
\label{sec:history-scale}

We next make the allocation depend on records. This experiment uses
fresh preparations, not the stationary trajectory of the preceding
section. At $L=64$ set $H_0=(\mu,\gamma)=(.1,0)$ with $g=0$.
Port $0$ reads the global mean, of variance $1/(L^2\mu)$.
Port $1$ reads one site, of variance
$v(\mu,\gamma)=L^{-2}\sum_k K_{\mu,\gamma}(k)^{-1}$.
Let
\[
 H_0^{\rm alt}=(.09,\gamma_\mu),\qquad
 v(.09,\gamma_\mu)=v(.1,0),\qquad
 H_1^{\rm alt}=(.1,1/12).
\]
The positive solution is $\gamma_\mu\simeq .0144867958104$;
monotonicity in $\gamma$ makes it unique. The equality, rather than
this rounded number, defines the preparation. Port $1$ is exactly
blind to $H_0^{\rm alt}$, and port $0$ to $H_1^{\rm alt}$.

In episode $t$, a hidden regime $J_t\in\{0,1\}$ selects the available
alternative. The source is $H_0$ or $H_{J_t}^{\rm alt}$ with equal
probability, independently of earlier source flags given the regime.
The regime starts uniform and changes with probability
$0<\varepsilon\leq1/2$ between episodes. Each reading uses a fresh
field and returns one bit:
$|Y|>\tau_b\sqrt{v_b(H_0)}$ at port $b$, with
$\tau_0=1.5222951844903065$ and $\tau_1=1.445716072545341$.
The baseline and informative-alternative probabilities are
\begin{equation}\label{eq:blind-probabilities}
 (p_0^0,p_0^1)\simeq(.1279351341,.1486894695),\qquad
 (p_1^0,p_1^1)\simeq(.1482568544,.1287055577).
\end{equation}
At the blind port the alternative probability is exactly the baseline.
The thresholds are specified designs, not claimed optimal quantizers.
All policies have $n=4000$ readings and preparations per episode and
choose the episode's split before collecting its records. The task is
to report whether that episode's source is the alternative.
Equal read counts do not equate the spatial support of a global mean
and a site measurement. This comparison stipulates equal read and
preparation prices; physical support and full apparatus energy remain
separate resources.

\begin{theorem}[Fixed-allocation bound and history-dependent allocation]
\label{thm:history-allocation}
Define the one-bit Bhattacharyya exponents
\[
 \beta_b=-\log\left(\sqrt{p_b^0p_b^1}
                  +\sqrt{(1-p_b^0)(1-p_b^1)}\right),\qquad
 \beta_{\max}=\max_b\beta_b.
\]
\begin{enumerate}
\item Every fixed split, or mixture of fixed splits independent of
the regime, has average error at least
\begin{equation}\label{eq:fixed-lower}
 B_n=\frac12\left(1-\sqrt{1-e^{-n\beta_{\max}}}\right),
\end{equation}
even if the decoder is told the current regime and keeps all past
records. For the stated experiment $B_{4000}\simeq .042754$.
\item A history policy with five retained bits has a finite Markov
description and a computable stationary error. One explicit policy
at $\varepsilon=10^{-4}$ has stationary error numerically evaluated
as $.034233$, below
\eqref{eq:fixed-lower}. Thus its advantage is over every fixed split,
not only a tested set of allocations.
\item At $\varepsilon=1/2$, history-dependent episode allocations
cannot beat the best history-free allocation under the same episode
protocol. This statement does not restrict policies which reallocate
after readings within the episode.
\end{enumerate}
\end{theorem}
\begin{proof}
For $m$ readings at an informative port, the two product laws have
Bhattacharyya coefficient $e^{-m\beta_b}$. Cauchy--Schwarz bounds
their total variation by $\sqrt{1-e^{-2m\beta_b}}$, so the equal-prior
Bayes error is at least
$b(m)=\frac12(1-\sqrt{1-e^{-2m\beta_{\max}}})$.
With the regime given, the other port carries no source information;
a split $(m,n-m)$ therefore has error at least
$[b(m)+b(n-m)]/2$. The decreasing function $b$ is convex on
$(0,\infty)$ and continuous at zero. Jensen's inequality gives
$b(n/2)=B_n$; independent randomization preserves the bound.
Past records cannot reveal the fresh source flag once $J_t$ is given.

For part (2), retain a believed regime $b\in\{0,1\}$ and a counter
$c\in\{1,\ldots,16\}$. Read only $b$. A coarse count at least $553$
or a fine count at most $553$ declares the alternative. A detection
raises $c$, capped at $16$. Without a detection, an independent fair
coin lowers $c$ with probability $1/2$; at zero switch $b$ and set
$c=1$. Nothing in this rule uses $J_t$ or the source flag. Its
transition matrix on $(J,b,c)$ has $64$ states. A binomial tail gives
each detection probability; the source flag is averaged with weight
$1/2$ and the regime changes after the update. Solving
$\pi T=\pi$, $\sum\pi=1$, and averaging the conditional task error
gives the stated numerical value. This is a deterministic evaluation
of a specified finite chain, not an adaptive optimality proof.

For part (3), $J_t$ is independent of all previous records and
private coins. Conditional on the history the allocation is fixed,
and current source data have their history-free law. The conditional
Bayes error cannot be below its minimum over fixed allocations.
Averaging proves the claim.
\end{proof}

The sharper minimum with the regime given is obtained by evaluating
$[e_0(m)+e_1(4000-m)]/2$ for all $4001$ splits, where
\[
 e_b(m)=\frac12\sum_{k=0}^{m}\min\{
 \operatorname{Bin}_{m,p_b^0}(k),
 \operatorname{Bin}_{m,p_b^1}(k)\}.
\]
It is $.095897$ at $m=1965$. The analytic bound, although less sharp,
already separates the history policy from the whole fixed class.
The same counter rule, with explicitly chosen cap $C$, yields:
\begin{center}\small
\begin{tabular}{rrr}
\toprule
$\varepsilon$ & $C$ & Stationary error\\
\midrule
$10^{-5}$ & $16$ & $.032754$\\
$10^{-4}$ & $16$ & $.034233$\\
$10^{-3}$ & $8$ & $.041516$\\
$10^{-2}$ & $4$ & $.081334$\\
\bottomrule
\end{tabular}
\end{center}
These are attainable policies, not minima over all histories or
universal switching thresholds. The unknown common gain problem of
Section~\ref{subsec:two-scales} has not been added to this experiment.

\begin{corollary}[A price interval for retained history]
\label{cor:history-price}
Give each correct source report value $V>0$ and define the exact
finite-comparison value
\[
 E_{\rm fix}=\min_{0\le m\le n}
               \frac{e_0(m)+e_1(n-m)}2.
\]
If the history policy's additional memory, randomization and update
charge per episode obeys
\[
 0\le\lambda<V(E_{\rm fix}-L_H),
\]
where $L_H$ is its stationary error, its net value exceeds that of
every fixed allocation, including mixtures, even with a regime-informed
decoder. The analytic bound $E_{\rm fix}\ge B_n$ gives the closed-form
subinterval $0\le\lambda<V(B_n-L_H)$.
At $\varepsilon=10^{-4}$, evaluation of the exact finite formulas gives
$E_{\rm fix}=.0958969785$, $L_H=.0342328069$ and
$(E_{\rm fix}-L_H)V\simeq.06166417V$, compared with the analytic
subinterval of approximately $.008521V$. Displayed rounding is not
an interval certificate.
\end{corollary}
\begin{proof}
At each fixed split the best regime-informed decoder has error
$[e_0(m)+e_1(n-m)]/2$: its other port is blind and the fresh source
flag is independent of past records conditional on the current regime.
A mixture of fixed splits cannot lower their minimum. Subtract the
common preparation and read charges; the gain in correct reports is
at least $V(E_{\rm fix}-L_H)$. Charging $\lambda$ gives the result.
\end{proof}

The public supplement \texttt{supplemental/history\_price} contains the
finite-state calculation and fixed-allocation enumeration for this corollary.
Here $\lambda$ is a task-value charge, not a thermodynamic estimate.
The retained automaton has $32$ states and needs five bits. Its
within-episode count has $4001$ states and needs twelve more transient
bits; acquisition timing, model storage, randomization, resetting and
computation remain in the apparatus bill. Silence is informative
because the contrast occurrence rate is specified: repeated absence
of an expected detection changes the future port choice. A rule
which switches only on a detection at the unread port never switches
when it takes no probe there. Such a frozen belief has long-run error
$\frac12 e_b(4000)+\frac14$, exceeding $.264$.
This is a failure of that update rule, not a claim that all
zero-probe policies fail. The constructed policy uses silence and
needs no second-port readings in these slowly changing regimes.

\section{A history chooses how to read and act on the same field}
\label{sec:field-control}

The preceding allocation experiment scores a report. Here a reading
is followed by an operation on its physical source. Take the stationary
process~\eqref{eq:OU}, with $K>0$ the finite real-space precision matrix.
Retain $R=a^TX_{-\Delta}+\xi_R$, where $a$ is a block vector. At
time zero either acquire one $Z_\ell=h_\ell^TX_0+\xi_\ell$ or take
no new reading. Independent Gaussian read noises have known positive
variances. These classical readings are assumed not to disturb $X$;
their apparatus and dissipation are not derived from the noise law.

After that decision, apply a constant force $Bu$ for $0<s<\Delta$,
with $B=a/\|a\|^2$ and $u\in\{-U,0,U\}$. Then
\begin{equation}\label{eq:field-force}
 dX_s=(-KX_s+Bu)\,ds+\sqrt2\,dW_s,\qquad
 Y_\Delta^u=Y_\Delta^0+\alpha u,\qquad
 \alpha=a^TK^{-1}(I-e^{-K\Delta})B>0,
\end{equation}
where $Y=a^TX$. The zero superscript is the unforced continuation of
the same field and bath noise, not a new preparation. In units with
bath temperature and mobility one, the force comes from the
nonnegative potential
$H_u(X)=\frac12(X-K^{-1}Bu)^TK(X-K^{-1}Bu)$.
The force pattern, finite amplitude, clock and read ports are inputs.
The task minimizes $\E[(Y_\Delta^u)^2+c_\ell]$, with zero read charge
when no read is acquired. These $c_\ell$ are task prices, not heat.

\begin{theorem}[Reading for a bounded physical action]
\label{thm:reading-action}
Suppose that, conditional on the authorized history $\mathcal H$,
the unforced target and candidate readings are jointly Gaussian.
Let $m=\E[Y_\Delta^0\mid\mathcal H]$ and put
\[
 q_\ell=
 \frac{\operatorname{Cov}(Y_\Delta^0,Z_\ell\mid\mathcal H)^2}
 {\Var(Z_\ell\mid\mathcal H)},\quad D=\alpha U,\quad
 \mathcal B(m)=2D(|m|-D/2)_+.
\]
For $N\sim N(0,1)$ define
\begin{equation}\label{eq:action-reading-value}
 G_\ell(m)=\E\mathcal B(m+\sqrt{q_\ell}N)-\mathcal B(m)-c_\ell.
\end{equation}
Choose the reading with the largest positive $G_\ell$, or no reading
if none is positive. Let $M$ be the revised conditional mean of the
unforced target. Apply zero force for $|M|\leq D/2$, and
$u=-U\operatorname{sign}M$ otherwise. The risk is
\begin{equation}\label{eq:field-control-risk}
 \E[(Y_\Delta^0)^2]-\E\mathcal B(m)
 -\E\max\{0,G_1(m),\ldots,G_p(m)\}.
\end{equation}
This minimizes the one-round objective over every measurable choice
using that history, at most one permitted new reading and the stated
three forces.
\end{theorem}
\begin{proof}
For mean $M$, the force changes conditional square loss by
$2\alpha uM+\alpha^2u^2$. Its minimum is $-\mathcal B(M)$.
Gaussian conditioning makes the revised mean, before receiving the
new result, distributed as $m+\sqrt{q_\ell}N$. The unforced target's
conditional second moment averages to its previous value. Thus
the reading's reduction is exactly $G_\ell$. Maximize pointwise
and average to obtain~\eqref{eq:field-control-risk}.
\end{proof}

Convexity of $\mathcal B$ makes the uncharged reading benefit
nonnegative. Its priced value depends on the history and on the
changes of optimal action at $M=\pm D/2$, not just on an information
rate. This is Bayesian experimental decision-making on a
scale-filtered carrier. Informative control is an established subject
\cite{KlenskeHennig}; the present Gaussian law and finite force menu
make its composition with imprint reading explicit.

\subsection{A finite-record case}

Set $L=32$, $\mu=.1$, $g=\gamma=0$, and $\Delta=.5$. The target and
old reading use a $4\times4$ block. Current ports use nested
$2\times2$ and $4\times4$ means. Every read noise variance is $.001$.
Choose $D=.15$, $c_2=.001$, $c_4=.014$. Finite lattice sums give
\begin{align*}
 v&=\Var(Y)=.1443687793,\qquad m=.6983747950R,\\
 \alpha&=.4012307531,\qquad U=.3738497083.
\end{align*}
and $q_2=.0079283594$, $q_4=.0272698232$.
The reader chooses scale $4$ for $|R|<.123905791$, scale $2$ for
$.123905791<|R|<.285137090$, and no read outside. Equalities may be
assigned either way. The thresholds are roots of the Gaussian value
functions, not fitted rules. With $d=D/2$,
\begin{equation}\label{eq:normal-force-benefit}
 \E\mathcal B(m+\tau N)=2D\sum_{\epsilon=\pm1}
 \left[\tau\phi\!\left(\frac{\epsilon m-d}{\tau}\right)
 +(\epsilon m-d)\Phi\!\left(\frac{\epsilon m-d}{\tau}\right)\right].
\end{equation}

\begin{table}[ht]
\centering\small
\begin{tabular}{lrrr}
\toprule
Next reading & Mean reads & Task risk & Switching work\\
\midrule
None & $0$ & $.100621$ & $.91202$\\
Always scale $2$ & $1$ & $.098280$ & $.94975$\\
Always scale $4$ & $1$ & $.103732$ & $1.02893$\\
History-selected & $.54545$ & $.097546$ & $1.02916$\\
\bottomrule
\end{tabular}
\caption{One continuation of the same field. Every row retains the old
reading and uses its optimal permitted force. Task risk includes the
read price; switching work is in bath-energy units and excludes the
reader and memory apparatus.}\label{tab:field-control}
\end{table}

Every choice independent of $R$ is a mixture of the first three rows,
even if its force decoder later uses $R$ and the new result. It cannot
beat their minimum. The history-selected rule does. Its material
square loss is $.093688$ and its mean read charge $.003858$. If all
information from $R$ is removed before the choice and action, the
best risk in this same menu is $.103835$. Neither comparison allows
an extra read or a field reset.

The strict advantage survives sixteen-bit records, finite coefficient
and threshold precision, and displacement error at most $10^{-5}$.
Appendix~\ref{app:force-precision} gives the implementation and register
account, certifying risk below $.097763$ against the ideal fixed-choice
minimum $.098280$. The analog interface, clock and physical storage
remain additional resources.

\subsection{Work and continued operation}

Switch $H_u$ on after reading and off at $\Delta$. The
force-dependent constant cancels between switches, giving the
pathwise mechanical work
\begin{equation}\label{eq:field-switch-work}
 W_{\rm sw}=uB^T(X_\Delta^u-X_0).
\end{equation}
This is the rectangular switching protocol in
Table~\ref{tab:field-control}, not $c_\ell$.
Finite ramps need not destroy the strict comparison. Replace the
force by $Buf(s)$, with $f$ rising linearly from zero to one in time
$\tau$, remaining one and falling to zero in the final $\tau$.
The response becomes
$\alpha_f=a^T\int_0^\Delta e^{-K(\Delta-s)}Bf(s)\,ds$.
Since $K>0$ and $\|a\|\|B\|=1$,
$|\alpha-\alpha_f|\leq2\tau$. Choosing
$\tau=5\times10^{-6}/(2U)$ bounds the displacement change by
$5\times10^{-6}$; the remaining $5\times10^{-6}$ may cover calibrated
amplitude error. This fits the finite-controller certificate.
The ramps have finite slope, though their timing and drive apparatus
remain supplied. Their actual work is
$-u\int_0^\Delta f'(s)B^TX_s\,ds$, obtained from the same potential;
the force-dependent offset integrates to zero.
The history-selected rule improves the declared objective but uses
more switching work than always reading scale $2$. It is not a
total-energy optimum. The measurement and memory parts of an
information-thermodynamic account require their own physical
implementation \cite{SagawaUeda2010}.

The resulting field can be read again. Known $K$, linear noisy
readings and retained controls give a Gaussian conditional-state
recursion on the same carrier \cite{ImprintChain}. After each force
its mean and covariance must be propagated; one cannot reapply the
old stationary prior or thresholds without that update. Reusing the
one-round rule is not a long-horizon optimality proof. Likewise, the
finite-record certificate above concerns this single continuation,
not an indefinitely running sixteen-bit controller.

\section{A control becomes a calibration reference on the same carrier}
\label{sec:active-reference}

The unknown-gain calculations used independent preparations. A control
can instead supply a reference on the field already being read, without
replacing its changed state by equilibrium. Choose one normalized real
spectral coordinate of the lattice, with known Laplacian eigenvalue
$s_0$, set $g=\gamma=0$, and let its unknown stable relaxation rate be
$\lambda=\mu+s_0>0$. A known force profile has projection $B>0$ on
this mode. On consecutive intervals it obeys
\[
 dX_t=(-\lambda X_t+Bu_i)\,dt+\sqrt2\,dW_t,\qquad
 Z_i=bX_{t_{i+1}}+\xi_i ,
\]
where $b>0$ is unknown and fixed, $\xi_i$ is independent centered
Gaussian read noise of known variance $r$. The mode, $B$, clock and
gain are the same for both reading durations. A force pattern or
gain allowed to change with the duration would define another problem.

Before each interval, use a fresh fair coin to choose $u_i=\pm U$,
$U>0$, independently of the entire past carrier, hidden parameters
and retained history. Alternate durations $\Delta$ and $2\Delta$,
with $\Delta>0$; neither the field nor the bank of old records is
reset. Put $f_\tau(\lambda)=(1-e^{-\lambda\tau})/\lambda$.

\begin{proposition}[Source and reader identified by paid responses]
\label{prop:active-reference}
For an interval of duration $\tau$, conditional on its full physical
past,
\[
 \E[u_iZ_i/U^2\mid\text{past}]
   =\kappa_\tau=bBf_\tau(\lambda).
\]
Thus the two response moments identify the relaxation rate and gain:
\begin{equation}\label{eq:active-identification}
 q=\frac{\kappa_{2\Delta}}{\kappa_\Delta}-1=e^{-\lambda\Delta},
 \qquad
 \lambda=-\frac{\log q}{\Delta},\quad
 \mu=\lambda-s_0,\quad
 b=\frac{\kappa_\Delta}{Bf_\Delta(\lambda)}.
\end{equation}
With known signed $B$, the response sign would also identify a
nonzero signed gain, unlike centered variance data. The finite
certificate below uses a positive gain interval. If $B$ is unknown,
only the product $bB$ is identified by these responses.

For a uniform finite-sample certificate, suppose
$\lambda\in[\lambda_-,\lambda_+]$, $b\in[b_-,b_+]$ with
$\lambda_-,b_->0$, and $\E X_0^2\le V_0<\infty$. Define
\[
 V_*=\max\left\{V_0,\frac1{\lambda_-}
                         +\left(\frac{BU}{\lambda_-}\right)^2\right\},
 \qquad
 Q=\frac{b_+^2V_*+b_+^2/\lambda_-+r}{U^2}.
\]
Collect $n$ intervals of each duration and average $u_iZ_i/U^2$
separately into $\widehat\kappa_\Delta,\widehat\kappa_{2\Delta}$.
For any $\eta>0$, both moment errors are at most $\eta$ except on
an event of probability at most $\min(1,2Q/(n\eta^2))$.
Let $\kappa_-=b_-Bf_\Delta(\lambda_+)$ and assume $\eta<\kappa_-$.
Clip the estimated $q$ to
$[e^{-\lambda_+\Delta},e^{-\lambda_-\Delta}]$ before taking its
logarithm. On the certified event,
\begin{align}
 |\widehat\lambda-\lambda|
 &\le e_\lambda:=
 \frac{3\eta}{(\kappa_--\eta)\Delta e^{-\lambda_+\Delta}},
 \label{eq:active-lambda}\\
 |\widehat b-b|
 &\le e_b:=\frac{\eta}{Bf_\Delta(\lambda_+)}
 +\frac{b_+\Delta^2}{2f_\Delta(\lambda_+)}e_\lambda .
 \label{eq:active-gain}
\end{align}
Clipping $\widehat b$ to $[b_-,b_+]$ preserves the latter bound.
\end{proposition}
\begin{proof}
The exact interval solution is
$X_{t+\tau}=e^{-\lambda\tau}X_t+Bf_\tau(\lambda)u_i+\zeta_i$,
with independent centered innovation of variance
$(1-e^{-2\lambda\tau})/\lambda$.
The fresh coin cancels the old-state term, even when that state is
nonzero or correlated with all old records. Multiplication by $u_i$
and conditional expectation give the response identity; the ratio
$f_{2\Delta}/f_\Delta=1+e^{-\lambda\Delta}$ proves identification.

The second-moment recursion is
\[
 \E X_{t+\tau}^2=e^{-2\lambda\tau}\E X_t^2+
 B^2U^2f_\tau(\lambda)^2+
 (1-e^{-2\lambda\tau})/\lambda .
\]
The chosen $V_*$ is invariant for either duration, since
$f_\tau(\lambda)^2/(1-e^{-2\lambda\tau})
=(1-e^{-\lambda\tau})/[\lambda^2(1+e^{-\lambda\tau})]
\le1/\lambda_-^2$.
The centered response samples are martingale differences with
conditional variance
\[
 U^{-2}\left[b^2e^{-2\lambda\tau}X_t^2+
 b^2(1-e^{-2\lambda\tau})/\lambda+r\right].
\]
Their unconditional second moments are at most $Q$. Orthogonality
gives each average variance at most $Q/n$, despite dependence of the
carrier across intervals. Chebyshev's inequality and a union bound
give the stated probability.
The true response ratio lies in $(1,2)$, so its error is at most
$3\eta/(\kappa_--\eta)$. Clipping $q$ cannot increase error, and
$\log$ is $e^{\lambda_+\Delta}$-Lipschitz on that interval.
Finally $|f_\Delta'|\le\Delta^2/2$ and
$f_\Delta\ge f_\Delta(\lambda_+)$ give~\eqref{eq:active-gain}.
\end{proof}

\begin{corollary}[Response calibration with unknown read noise]
\label{cor:active-noise}
In Proposition~\ref{prop:active-reference}, replace the Gaussian
read errors by mutually independent centered errors, independent of
the carrier past, bath increments and force coins, with
$\E\xi_i^2\le r_+<\infty$. Their distributions and actual variances
may be unknown and may differ between readings. The response identity,
source--gain identification and finite parameter-error certificate
remain valid with $r_+$ in place of $r$ in $Q$.
\end{corollary}
\begin{proof}
The read error contributes zero to the conditional response mean.
Independence removes its cross terms with the old-state and bath
terms. Its contribution to the conditional variance is
$\E\xi_i^2/U^2\le r_+/U^2$.
The carrier's second-moment recursion is unchanged because read errors
do not enter its force coin or dynamics. The centered responses are
therefore still martingale differences with the same upper bound $Q$
after the substitution. Orthogonality, Chebyshev's inequality and the
deterministic ratio and logarithm bounds prove the claim.
\end{proof}

This extension does not identify the read-noise law. Nor does it extend
the Gaussian likelihoods, entropy rates or conditional-state filters
used elsewhere in the paper. With unknown parameters, or non-Gaussian
read noise, those Gaussian state updates cannot be recovered by
inserting fitted parameters and a variance cap. The response added by
the next force, however, depends on the same linear dynamics and
conditional zero mean, so Corollary~\ref{cor:active-next} also uses this
certificate. Without a variance cap the population response identity
can still hold, but the uniform finite-budget bound above is unavailable.

\begin{corollary}[Finite records for response calibration]
\label{cor:active-quantization}
In Corollary~\ref{cor:active-noise}, clip each endpoint record to
$[-A,A]$ and round to $2^{B_{\rm rec}}$ equally spaced levels, with
step $\delta=2A/(2^{B_{\rm rec}}-1)$. Put
$M_Z=b_+^2V_*+r_+$ and
$e_{\rm rec}=[\delta/2+M_Z/(4A)]/U$.
For any $\eta,\epsilon>0$ with
$\eta+\epsilon<\kappa_-$, the parameter and next-action bounds
remain valid with $\eta+\epsilon$ in place of $\eta$, except on an
event of probability at most
\[
 \min\left\{1,\frac{2Q}{n\eta^2}
                 +\frac{2e_{\rm rec}}{\epsilon}\right\}.
\]
Thus a sufficiently large but finite range and record register
retain every fixed positive-tolerance certificate.
\end{corollary}
The proof is in Appendix~\ref{app:response-record-proof}.
The accumulation can be exact on integer level indices: a signed sum
for each duration needs at most
$B_{\rm rec}+\lceil\log_2(n+1)\rceil+2$ bits, with the record range,
force scale and subsequent parameter arithmetic supplied separately.
Quantization does not manufacture an analog measuring device or
a cost-free cyclic processor.

No independent-read variance formula was used. The fresh randomness
creates orthogonal response errors, not independent field samples.
The price is $2n$ actual force intervals and endpoint readings,
$3n\Delta$ elapsed time, the specified spatial force/read support,
coins, moment accumulation and retained records. Corollary~\ref{cor:active-quantization} gives the additional finite-record
error and the accumulation storage needed for this certificate. For each fixed positive gain one can use its
actual finite constants. Without a positive lower gain bound this is
not a uniform fixed-budget claim as $b\to0$.
For failure probability at most $\delta_{\rm fail}$, the sufficient
count is $n\ge2Q/(\delta_{\rm fail}\eta^2)$. This Chebyshev budget
is conservative, not an efficiency or optimal-design claim.

\begin{corollary}[Calibration selects the next time scale and action]
\label{cor:active-next}
Permit force amplitudes in $[0,U_{\max}]$, rather than only the
three-action menu of Section~\ref{sec:field-control}.
Let $q_*\in(0,1)$ be a desired response decay. After the calibration
history, set $\widehat\tau=-\log(q_*)/\widehat\lambda$.
On the preceding certified event,
\[
 |\widehat\tau-\tau_*|
 \le\frac{-\log(q_*)}{\lambda_-^2}e_\lambda,
 \qquad \tau_*=-\log(q_*)/\lambda .
\]
For a desired readout displacement
$0\le D_*\le U_{\max}b_-B(1-q_*)/\lambda_+$, install the legal force
\[
 \widehat u=\frac{D_*\,\widehat\lambda}
                  {\widehat b B(1-q_*)}\le U_{\max}
\]
for this next interval. With $\tau_{\max}=-\log(q_*)/\lambda_-$,
its actual mean displacement relative to the unforced continuation
of the same current state differs from $D_*$ by at most
\[
 U_{\max}B\left[\tau_{\max}e_b+
                       b_+\tau_{\max}^2 e_\lambda/2\right].
\]
\end{corollary}
\begin{proof}
The reciprocal function is $\lambda_-^{-2}$-Lipschitz.
The predicted displacement is exactly
$\widehat b B f_{\widehat\tau}(\widehat\lambda)\widehat u=D_*$.
At that same actual duration,
$|bBf_{\widehat\tau}(\lambda)-
 \widehat b Bf_{\widehat\tau}(\widehat\lambda)|
\le B[\tau_{\max}e_b+b_+\tau_{\max}^2e_\lambda/2]$.
Multiply by the legal amplitude.
\end{proof}

The comparison concerns the response added to a current, generally
nonstationary field. It does not erase the calibration's perturbation.
The old state contributes to the absolute next record and must be
propagated. Once unknown parameters are inferred from these records,
a posterior over parameters and state is generally a mixture, not
one Gaussian state. The known-parameter Gaussian rule of
Theorem~\ref{thm:reading-action} may not be applied to that mixture by
silently replacing it with a fitted equilibrium law.

Force ramps also belong to the model. A normalized taper replaces
$f_\tau$ by
$\int_0^\tau e^{-\lambda(\tau-s)}f(s)\,ds$; arbitrary tapers do not
preserve~\eqref{eq:active-identification}. The exact identity above
uses the stated rectangular force. Its potential is nonnegative,
$H_u(X)=\lambda(X-Bu/\lambda)^2/2$, and a switch from $u$ to $v$
costs the actual potential difference at the current $X$, not a
price attached to an observation. Finite amplitudes and durations
do not by themselves supply a heat, memory-reset or complete-apparatus
work ledger.

\section{From trace-carrying to organized reading}

A source contrast, its carrier and the available reading are distinct.
The zero mode can miss a contrast present in the field; one spectral
band can confuse a source change with reader gain. Temporal filtering
can instead make past records useful when the coarse present is
insufficient. This supplies a precise retention-and-reuse interface
to Stable Hulls~\cite{StableHulls}, with a source-specific meaning of
access dependence~\cite{HAFFA}.

Records also change the next physical operation. The regime protocol
selects a port, the Gaussian continuation selects a reading and force,
and paid responses calibrate a future duration and command on a
changed carrier. Their preparation and noise laws remain different;
unknown gain, reader drift and hidden switching cannot be combined by
identifying their estimates. The Imprint Chain supplies separate
use and maintenance permissions~\cite{ImprintChain}; this field does
not derive its body's repair law. Port formation, control access and
first-person experience remain separate explanatory questions.

\section{Conclusion}

Scale-stable readability is not microscopic reconstruction.
Block averaging in the specified Gaussian field retains substantial
mass information while suppressing higher-gradient contrasts.
The same operation leaves a dynamic consequence: past and future are
conditionally correlated when the coarse present is insufficient.
A finite retained record exploits that correlation, with strictly
better risk than any current-only predictor. Additional spectral
access separates a source parameter from the reader's own gain.
Independent block readings supply an exact finite-sample calibration
law. Noisy readout distinguishes two further questions: whether a
source is identifiable at all, and whether the available probe budget
can identify it reliably. Known probe variations can supply a
calibration reference even at one spatial scale.

Filtering now has two consequences for a reader. It gives retained
history a prediction value, and it changes which scale is worth
reading when the source or reader changes. A finite history can govern
the next allocation and beat every fixed allocation under an explicit
price interval. On a single continuing field, it can also choose a
noisy block reading and a bounded physical force. The optimal
one-round rule, its finite-record advantage and its switching work
make prediction, inquiry and actual change parts of the same
experiment. The active response reference adds a further step:
a history can learn which source and reader parameters govern its
actual intervention, then select a new time scale and command on
the changed carrier. The moment certificate prices a finite experiment
but does not optimize it. The formation of ports, known force
relations and maintenance capacities remains a physical question.
Unknown centered independent read noise with a known variance cap
preserves this response route, without turning the other Gaussian
inference protocols into noise-independent laws.

\appendix
\section{Numerical evaluation}

Spectral sums use the complete finite grid with analytic derivatives.
Chernoff optimization is one-dimensional on $[0,1]$. For stability its
per-mode summand is evaluated as
$\frac12\{\log[1+u(e^z-1)]-uz\}$, where $z=\log(v_1/v_0)$,
using \texttt{log1p} and \texttt{expm1}. Finite bins are integrated
before computing their discrete Chernoff coefficient.

For an $8\times8$ lattice, direct physical-space covariance and block
matrices agree with \eqref{eq:block} to $3.6\times10^{-15}$.
A sample of $6000$ fields at $L=32,\ell=4$ has mean
measured-to-predicted spectral variance ratio $.995762$, without
removing an overall normalization. Direct Gaussian regression agrees
with \eqref{eq:risks} to about $10^{-17}$.

For Table~\ref{tab:ratio-counts}, scalar block variances are also
evaluated directly on the fine lattice as
\[
 \sigma_\ell^2(H)=\frac1{L^2\ell^4}
 \sum_k\frac{|W_\ell(k_1)W_\ell(k_2)|^2}{K_H(k)}.
\]
The equal-count geometric decision boundaries give success
probabilities from $F$ cumulative and survival functions. Searching
integer counts and checking each preceding count locates the listed
thresholds; no density-grid quadrature is used.

Table~\ref{tab:self-scale} uses the finite Fourier transform of each
physical block and the spectral density~\eqref{eq:reading-density}.
Periodic midpoint grids of $8192$ and $16384$ frequencies agree within
$9\times10^{-16}$ for these listed rates. This agreement is a numerical
convergence check, not an interval certificate for arbitrary spectra.
The history policies use analytic binomial tails, an enumerated
finite transition matrix and its normalized stationary solution;
the stationarity residual is below $3\times10^{-16}$. The fixed
regime-informed comparison enumerates all integer splits. Its
likelihood-threshold tails agree with direct binomial sums through
$m=32$ within $9\times10^{-16}$. No Monte Carlo estimate or policy
search is needed for the stated constructed-policy comparisons.

For the active response reference, conditional means were checked
by explicit averaging over the two forces at four nonzero or zero
old-state values. A diagnostic uses $64$ continuing carriers, each
with $20{,}000$ intervals at each duration, starting at $X_0=2$.
At $(\lambda,b,B,U,\Delta,r)=(.7,1.3,.8,.6,.5,.01)$, population
response moments are $.4387491$ and $.7479304$; empirical replica
means are $.4416915$ and $.7465218$. Their comparison uses the
martingale variances accumulated from the actual second-moment
recursion, not an independent-field approximation.
Perturbed response moments at the four corners of a fixed error
square satisfy the parameter and next-action bounds. A force sign
aligned with the old state instead gives a nonzero response bias,
confirming the independence requirement. These checks do not
certify force apparatus heat or optimize the conservative budget.

\subsection{Finite response records}
\label{app:response-record-proof}

\begin{proof}[Proof of Corollary~\ref{cor:active-quantization}]
For the quantizer $\mathsf q$, the pathwise error is bounded by
$\delta/2+(|Z|-A)_+$.
The inequality $(|z|-A)_+\le z^2/(4A)$ and
$\E Z_i^2\le M_Z$ give $\E|\mathsf q(Z_i)-Z_i|
\le\delta/2+M_Z/(4A)$.
Each duration's response average therefore has mean absolute
quantization error at most $e_{\rm rec}$.
Markov's inequality and a union bound put both such errors below
$\epsilon$ with failure at most $2e_{\rm rec}/\epsilon$.
Combine this with the unquantized martingale event. The remaining
ratio, clipping and action bounds are deterministic.
Taking $A$ large and then $B_{\rm rec}$ large makes $e_{\rm rec}$
arbitrarily small; no Gaussian tail law was used.
\end{proof}

\section{Finite records for the reading-and-force rule}
\label{app:force-precision}

\subsection{Rational arithmetic for the prediction rule}

An entirely finite-arithmetic variant uses rational $A=3$, the same
$10$-bit records, and rounds $a_1,a_2$ and the output to the nearest
$2^{-16}$. Intermediate rational arithmetic is finite.
Coefficient and output rounding add at most $(2A+1)2^{-17}$ to the
predictor error. Adding that amount to $(|a_1|+|a_2|)\eta$ before
squaring in \eqref{eq:finite} gives risk at most $.021433187<R_1$.
The parameters are known and the extra registers must be counted.
The advantage therefore does not require an infinite-precision
archive, though it remains relative to a specified model and interface.

\subsection{The reading-and-force controller}

Both old and newly acquired scalar records use the nearest-point uniform
quantizer on $[-4,4]$ with $2^{16}$ levels, including the endpoints, and
saturation outside the interval. Decoder coefficients are rounded to
sixteen binary fractional bits; the two $R$-thresholds to twenty, giving
$\epsilon_R\le 2^{-21}$. This is distinct from the ten-bit prediction
variant above.
There are sixteen retained data bits and sixteen transient data bits;
five fixed decoder coefficients use eighty fractional bits in total
and the two thresholds forty. A forty-bit signed accumulator and a
twenty-one-bit revised-mean register suffice for this arithmetic
construction. Program, acquisition clock, analog interface and
physical storage remain additional resources; these are sufficient
register sizes, not minimal ones.


\paragraph{Quantization and decision regret.}
For a centered Gaussian record of variance $V$, clipping and uniform
quantization on $[-A,A]$ have $L^1$ error at most
\[
 \eta(V)=\frac{A}{2^B-1}
 +\sqrt{2V[z\phi(z)+\overline\Phi(z)]},\qquad z=A/\sqrt V.
\]
Use the $L^2$ clipping bound of Corollary~\ref{cor:finite} and
Cauchy--Schwarz. Write the old mean as $\beta R$ and the revised
mean as $b_\ell R+d_\ell Z_\ell$, where
\[
 d_\ell=\frac{\operatorname{Cov}(Y_\Delta^0,Z_\ell\mid R)}
 {\Var(Z_\ell\mid R)},\qquad
 b_\ell=\beta-d_\ell\frac{\operatorname{Cov}(Z_\ell,R)}{\Var R}.
\]
Include $b_0=\beta$ for no reading, and set $C_R=\max_{\ell\geq0}|b_\ell|$.

The function $\mathcal B$ is $2D$-Lipschitz, and each $G_\ell$ is
$4D$-Lipschitz. Maximizing scores at a perturbed old mean costs at
most twice their largest perturbation. Choosing the force at a
perturbed revised mean costs at most $4D$ times its error.
The record-induced regret is therefore bounded by
\[
 8D|\beta|\eta_R+
 4D\left(C_R\eta_R+\sum_\ell|d_\ell|\eta_\ell\right).
\]
The sum bounds errors at a selected port using its unconditional
Gaussian marginal. It does not grant unchosen readings to the policy.

In the example, both innovation standard deviations exceed $d=D/2$.
The larger-variance minus smaller-variance smoothed benefit decreases
with $|m|$. Each reading gain over no reading increases up to $d$
and decreases afterwards. Together with the signs at zero and the
crossing, these give precisely the three selection intervals.
To verify the difference's monotonicity, differentiate
\eqref{eq:normal-force-benefit}. For $0<m<d$ and $\tau\geq d$,
\[
 \log\frac{d+m}{d-m}\geq\frac{2m}{d}
 \geq\frac{2dm}{\tau^2};
\]
the derivative of the smoothed benefit in $m$ consequently decreases
with $\tau$. For $m\geq d$ the same sign follows directly from the
two densities. Rounding the $R$ boundaries with error $\epsilon_R$
adds at most $8D|\beta|\epsilon_R$: their competing scores agree
at an exact boundary within that distance.

Coefficient and output rounding on a grid of sixteen binary
fractional bits contributes mean error at most $(2A+1)2^{-17}$,
and hence regret at most $4D(2A+1)2^{-17}$.
An actual force displacement error $\epsilon_D$ contributes at most
$2\epsilon_D(\sqrt v+D)+\epsilon_D^2$, because the unforced target
has second moment $v$ and the applied displacement is bounded by $D$.
With $B=16$, $A=4$, $\epsilon_R=2^{-21}$ and
$\epsilon_D=10^{-5}$, total regret is at most $.000164788$.
Numerical evaluation gives risk below $.097711$.
A conservative directed-interval calculation also certifies the strict
comparison. Put $s=\sqrt{\Var(m)}$ and $\tau_{\max}=\max_\ell\sqrt{q_\ell}$.
On the positive half-line the product of the folded Gaussian density
and $\max(0,G_2,G_4)$ is Lipschitz with constant at most
\[
 \mathcal L=
 \frac{4D\sqrt{2/\pi}}{s}
 +\frac{2D\tau_{\max}\sqrt{2/\pi}}{s^2}.
\]
Use $1024$ midpoint intervals on $[0,.3]$ and subtract
$\mathcal L(.3)^2/(4\cdot1024)$ from the gain integral estimate.
Dropping its nonnegative remainder beyond $.3$ preserves a lower
bound. Evaluate the finite lattice sums and interval arithmetic at
55 decimal digits. The Gaussian CDF is evaluated with the integrated
exponential Taylor polynomial of order $48$, using remainder at most
$|x|^{99}/(2^{49}49!\,99\sqrt{2\pi})$.
For the clipping bound one may use
$\overline\Phi(z)\leq\phi(z)/z$.
This gives finite-controller risk below $.097763$ and fixed-choice
risk above $.09827985$. The certificate concerns this one-round
comparison, not an infinite-duration apparatus or formal verification
of the interval software.

\begin{thebibliography}{99}
\bibitem{KlenskeHennig}
E.~D.~Klenske and P.~Hennig,
\emph{Dual Control for Approximate Bayesian Reinforcement Learning},
Journal of Machine Learning Research \textbf{17}(127), 1--30 (2016),
\href{https://jmlr.org/papers/v17/15-162.html}{journal text}.
\bibitem{SagawaUeda2010}
T.~Sagawa and M.~Ueda,
\emph{Generalized Jarzynski Equality under Nonequilibrium Feedback Control},
Physical Review Letters \textbf{104}, 090602 (2010),
\href{https://arxiv.org/abs/0907.4914}{arXiv:0907.4914}.
\raggedright
\bibitem{Gray2006}
R.~M.~Gray, \emph{Toeplitz and Circulant Matrices: A Review},
Foundations and Trends in Communications and Information Theory
\textbf{2}(3), 155--239 (2006),
\href{https://ee.stanford.edu/~gray/CIT006-journal.pdf}{author's text}.
\bibitem{ImprintChain}
S.~Liu, \emph{The Imprint Chain: Records, Self-Calibration, and
Physical Feedback}, 2026;
series record DOI: 10.5281/zenodo.20151594.
\bibitem{NISTF}
NIST/SEMATECH, \emph{e-Handbook of Statistical Methods},
Section~1.3.6.6.5, ``F Distribution,''
\href{https://www.itl.nist.gov/div898/handbook/eda/section3/eda3665.htm}{online reference}.
\bibitem{WorldImprint}
S.~Liu, \emph{The World as Imprint: From Trace-Carrying Structure
to Self-Reading Mind}, Zenodo (2026),
\href{https://doi.org/10.5281/zenodo.20122718}{doi:10.5281/zenodo.20122718}.
\bibitem{Machta2013}
B.~B.~Machta, R.~Chachra, M.~K.~Transtrum and J.~P.~Sethna,
\emph{Parameter space compression underlies emergent theories and
predictive models}, Science \textbf{342}, 604--607 (2013),
\href{https://arxiv.org/abs/1303.6738}{arXiv:1303.6738}.
\bibitem{Raju2018}
A.~Raju, B.~B.~Machta and J.~P.~Sethna,
\emph{Information loss under coarse graining: A geometric approach},
Physical Review E \textbf{98}, 052112 (2018),
\href{https://doi.org/10.1103/PhysRevE.98.052112}{doi:10.1103/PhysRevE.98.052112}.
\bibitem{Mori1965}
H.~Mori, \emph{Transport, collective motion, and Brownian motion},
Progress of Theoretical Physics \textbf{33}, 423--455 (1965),
\href{https://doi.org/10.1143/PTP.33.423}{doi:10.1143/PTP.33.423}.
\bibitem{StableHulls}
S.~Liu, \emph{Stable Hulls as Imprint Integration: Relations, Anchors
and the Cost of Joint Reading}, 2026,
\href{https://doi.org/10.5281/zenodo.20151508}{doi:10.5281/zenodo.20151508}.
\bibitem{HAFFA}
S.~Liu, \emph{Physical Access and Correlation Geometry: The Holographic Alaya--Field Framework}, Zenodo (2026),
\href{https://doi.org/10.5281/zenodo.18361706}{doi:10.5281/zenodo.18361706}.
\bibitem{TDIII}
S.~Liu, \emph{Keeping a Self-Model Calibrated: Drift Detection, the
Price of Keeping Up, and Reliability in a Closed Loop} (T-DOME III),
2026.
\bibitem{FoundationIII}
S.~Liu, \emph{Foundation III: The Perfuming Loop.
Held Content, Learned Frames and Physical Return}, 2026,
section on finite feedback; series record DOI: 10.5281/zenodo.19642189.
\end{thebibliography}
\end{document}
